\documentclass[11pt,a4paper]{article}
\usepackage[utf8]{inputenc}
\usepackage[T1]{fontenc}
\usepackage{lmodern}
\usepackage[margin=2.2cm]{geometry}
\usepackage{amsmath,amssymb,amsfonts}
\usepackage{graphicx}
\usepackage{booktabs}
\usepackage{float}
\usepackage{hyperref}
\usepackage{cleveref}
\usepackage{caption,subcaption}
\usepackage{xcolor}
\usepackage{multirow}
\usepackage{natbib}
\usepackage{tcolorbox}
\hypersetup{colorlinks=true,linkcolor=blue!60!black,citecolor=blue!60!black,urlcolor=blue!60!black}
\graphicspath{{../results/figures/}}

\newtcolorbox{intuition}{
  colback=blue!5!white,
  colframe=blue!50!black,
  fonttitle=\bfseries,
  title=Clinical Intuition,
  boxrule=0.5pt,
  arc=2pt,
  left=6pt, right=6pt, top=4pt, bottom=4pt
}

\title{\textbf{The Sim-to-Real Uncertainty Gap in Quantitative MRI:\\Characterization, Benchmark, and Counterfactual Correction}}
\author{Sreenath Kyathanahally}
\date{June 2026}

\begin{document}
\maketitle

\begin{abstract}
Machine learning models trained on simulated MRI data are increasingly used to extract quantitative measurements from real brain scans. But can these models tell us when they might be wrong? We provide evidence that they usually cannot. When tested on real patient data without any fine-tuning, the models still produce reasonable measurements, but their confidence estimates break down---they can no longer distinguish reliable predictions from unreliable ones. We call this the \textbf{sim-to-real uncertainty gap}. We demonstrate that this gap can be fixed with a quick calibration step using only 5\% of real scan data. To help the community study and solve this problem, we release \textbf{qMR-FailureBench}, a standardized benchmark of 60,000 simulated MRI signals with five evaluation tasks. We also show that our system can not only detect unreliable measurements, but identify \emph{why} they failed and attempt to correct them, reducing errors by 39.6\%.
\end{abstract}

\section{Introduction}
\label{sec:intro}

MRI scanners are one of the most important tools in modern medicine. Unlike X-rays or CT scans, MRI uses magnetic fields and radio waves---not radiation---to create detailed images of the body's internal structures. While most people know MRI for its pictures, there is a growing push to extract \emph{quantitative measurements} from MRI scans---numbers that describe tissue properties like how quickly water molecules relax back to equilibrium. These measurements, called $T_1$ and $T_2$ relaxation times, can reveal early signs of diseases like multiple sclerosis, Alzheimer's, or cancer, often before symptoms appear.

\begin{intuition}
\textbf{Conventional MRI} produces weighted images (``this region appears bright or dark depending on acquisition parameters''). \textbf{Quantitative MRI (qMRI)} produces tissue-specific biophysical parameters (``this voxel has $T_1 = 800$~ms, consistent with normal-appearing white matter''). These numbers enable objective, cross-scanner, cross-site comparisons---functioning as imaging-derived biomarkers analogous to serum laboratory values.
\end{intuition}

Extracting these numbers from raw MRI signals is mathematically complex. In recent years, researchers have turned to deep learning: training neural networks on millions of simulated (synthetic) MRI signals to learn the mapping from raw signal to tissue parameters. This approach works remarkably well in the lab, but there is a critical open question:

\begin{center}
\textit{When the model makes a mistake on a real patient scan, does it know it is wrong?}
\end{center}

A model that produces a $T_1$ value of 1200~ms when the true value is 800~ms is dangerous if it also reports high confidence. In clinical practice, such ``confidently wrong'' measurements could lead to misdiagnosis.

\begin{intuition}
Consider a radiologist who has trained exclusively on textbook cases from a single hospital. When transferred to a different hospital with different scanner hardware, patient demographics, and imaging protocols, the radiologist may still render diagnoses---but their calibrated sense of ``how certain am I?'' breaks down. They have lost the internal calibration that comes from experience with the local data distribution. Our models exhibit the same failure mode when moving from synthetic training data to real scanner data.
\end{intuition}

\textbf{Contributions.}
(1)~We provide evidence for a \textbf{sim-to-real uncertainty gap} in qMRI and show it is fixable with 5\% of real data.
(2)~We release \textbf{qMR-FailureBench} (60K signals, 5 tasks, \url{https://doi.org/10.5281/zenodo.20632105}) as a standardized evaluation framework.
(3)~We build a system that \textbf{detects, explains, and corrects} failures, reducing errors by 39.6\%.
(4)~We observe this gap in both MRF and a preliminary MRS study, suggesting a systemic property of physics-informed evidential networks.

\section{Background and Related Work}
\label{sec:related}

\paragraph{How MRI scanners can fail.}
Three types of physical imperfections commonly corrupt MRI signals:
\begin{itemize}
  \item \textbf{$B_0$ off-resonance:} The main magnetic field is never perfectly uniform. Spatial inhomogeneity shifts the local Larmor frequency, causing phase accrual across the echo train. In MRF, this manifests as a systematic shift in the apparent $T_1$/$T_2$ dictionary match.
  \item \textbf{$B_1^+$ transmit inhomogeneity:} The radio-frequency excitation field varies spatially across the imaging volume. Flip angles deviate from their prescribed values, altering the nonlinear Bloch equation trajectory and biasing parameter estimates.
  \item \textbf{Patient motion:} Bulk head movement during acquisition introduces k-space phase errors that manifest as ghosting, blurring, and signal loss---particularly damaging in multi-echo or fingerprinting sequences.
\end{itemize}

In practice, these artifacts \emph{co-occur}---we call this \textbf{entanglement}. A patient who moves during a scan simultaneously introduces motion artifacts \emph{and} perturbs the local $B_0$ field through susceptibility changes.

\begin{intuition}
In a typical clinical MRF acquisition, $B_0$ inhomogeneity near the air-tissue interface at the skull base, $B_1^+$ shading from the transmit coil geometry, and involuntary head motion are not independent events---they are concurrent physical processes. A correction algorithm that addresses only one source of bias while ignoring the others will propagate residual errors into the parameter map. Our \texttt{PhysicsCorruptor} is designed to simulate exactly this entangled regime.
\end{intuition}

\paragraph{Uncertainty in deep learning.}
Standard neural networks output a single number with no indication of reliability. Deep Ensembles~\citep{lakshminarayanan2017ensembles} train $M$ models and measure disagreement; MC-Dropout~\citep{gal2016dropout} approximates Bayesian inference via stochastic forward passes. Both require $M \gg 1$ forward passes at inference time. \textbf{Evidential Deep Learning}~\citep{sensoy2018evidential,amini2020deep} places a prior over the likelihood parameters, enabling uncertainty decomposition in a single forward pass---critical for real-time clinical deployment.

\paragraph{Sim-to-real transfer.}
The gap between synthetic training and real deployment is well-documented in computer vision and robotics~\citep{tobin2017domain}. In medical imaging, similar challenges arise in synthetic CT generation~\citep{han2019deep}. However, whether \emph{uncertainty calibration} survives this transfer has not been systematically studied for qMRI.

\section{Methods}
\label{sec:methods}

\subsection{qMR-FailureBench}

We created qMR-FailureBench to give the community a standardized testbed for studying qMRI failures:

\begin{table}[H]
  \centering
  \caption{qMR-FailureBench evaluation tasks.}
  \label{tab:tasks}
  \begin{tabular}{@{}lll@{}}
    \toprule
    \textbf{Task} & \textbf{Metric} & \textbf{Clinical Question} \\
    \midrule
    Failure detection & AUROC & ``Is this measurement trustworthy?'' \\
    Corruption attribution & F1 & ``\emph{Why} did it fail?'' \\
    Severity estimation & MAE & ``How bad is the artifact?'' \\
    Counterfactual repair & $\Delta$MAE & ``Can we fix it?'' \\
    Sim-to-real calibration & Spearman $\rho$ & ``Does this work on real scanners?'' \\
    \bottomrule
  \end{tabular}
\end{table}

\subsection{Entangled Physics Corruption}

Our \texttt{PhysicsCorruptor} injects realistic scanner imperfections into simulated signals. For each sample, independent Bernoulli trials (controlled by probability weights $w_{B_0}, w_{B_1}, w_{\text{motion}}$) determine which corruption sources are active, with the constraint that at least one is always applied. This models the clinical reality that artifacts rarely occur in isolation.

\begin{table}[H]
  \centering
  \caption{Corruption parameter ranges and distributions.}
  \label{tab:corruption_params}
  \begin{tabular}{@{}llll@{}}
    \toprule
    \textbf{Corruption} & \textbf{Parameter} & \textbf{Range} & \textbf{Distribution} \\
    \midrule
    $B_0$ off-resonance & $\Delta f$ (Hz) & $[-80, 80]$ & Uniform \\
    $B_1^+$ scaling & $\lambda$ & $[0.6, 1.4]$ & Uniform \\
    Motion translation & $\delta$ (voxels) & $[-8, 8]$ & Uniform integer \\
    Motion rotation & $\theta$ (degrees) & $[-15, 15]$ & Uniform \\
    \bottomrule
  \end{tabular}
\end{table}

Our training dataset comprises 49,950 MRF signals (1,000 timepoints each, 80/20 train/validation split) and 10,000 MRS spectra (2,048 points each), generated across 3 vendors (Siemens, Philips, GE) and 3 field strengths (1.5T, 3T, 7T).

\subsection{Dual-Head Architecture}

Our network has two heads sharing one backbone:
\begin{enumerate}
  \item \textbf{Regression head:} Outputs a Normal-Inverse-Gamma distribution $(\gamma, \nu, \alpha, \beta)$ per tissue parameter. The parameter $\gamma$ is the predicted mean; $\nu$ and $\alpha$ encode epistemic and aleatoric evidence respectively; $\beta$ is the scale parameter.
  \item \textbf{Severity head:} Estimates corruption magnitudes ($\hat{\Delta f}$ in Hz, $\hat{\lambda}$ dimensionless, $\hat{\delta}$ in voxels).
\end{enumerate}

\subsection{Evidential Regularizer}

Standard NIG training can minimise the NLL by increasing evidence ($\nu, \alpha$) without improving predictions. Our regularizer couples evidence to accuracy:
\begin{equation}
  \mathcal{L}_{\text{ER}} = |y - \gamma| \cdot (2\nu + \alpha)
\end{equation}
This implements a soft constraint: high evidence is only rewarded when predictions are accurate. When the model encounters out-of-distribution artifacts, the regularizer forces evidence to decrease, inflating epistemic uncertainty.

\begin{intuition}
In a well-calibrated diagnostic system, confidence should be proportional to accuracy. A radiologist who is ``very certain'' about 90\% of their diagnoses and ``uncertain'' about the difficult cases is well-calibrated. One who is ``very certain'' about everything---including the cases they get wrong---is dangerous. The Evidential Regularizer enforces this calibration principle mathematically.
\end{intuition}

\subsection{Counterfactual Correction Pipeline}

Our pipeline: \textbf{Detect} (epistemic $> \tau$) $\to$ \textbf{Attribute} (severity head estimates $\hat{\Delta f}, \hat{\lambda}, \hat{\delta}$) $\to$ \textbf{Correct} (apply inverse corruption: scale by $1/\hat{\lambda}$, phase-shift by $-\hat{\Delta f}$, translate by $-\hat{\delta}$) $\to$ \textbf{Re-predict}.

\begin{intuition}
In MRI artifact correction, the standard workflow is: acquire data, reconstruct, identify artifact, apply correction. Our system automates the ``identify'' and ``correct'' steps by learning to estimate the corruption parameters directly from the signal, then inverting the estimated corruption in k-space or signal domain before re-prediction.
\end{intuition}

\section{Results}
\label{sec:results}

\subsection{The Sim-to-Real Uncertainty Gap}

Our central finding: when a model trained on synthetic data is tested on real in-vivo brain scans (qMRLab~\citep{qmrlab}, 4,668 voxels), regression transfers (MAE = 64.7~ms) but uncertainty calibration collapses.

\begin{intuition}
The model produces reasonable $T_1$ estimates on real brain data, but its uncertainty scores become uninformative---it cannot distinguish a clean white-matter voxel from one near a susceptibility boundary. This is the epistemic equivalent of a radiologist who has lost the ability to flag cases for peer review: the diagnoses may still be mostly correct, but the quality-control mechanism has broken down.
\end{intuition}

\begin{figure}[H]
  \centering
  \includegraphics[width=\textwidth]{sim_to_real_gap.png}
  \caption{The sim-to-real uncertainty gap on qMRLab brain data. Regression transfers but uncertainty does not. Isotonic repair recovers calibration.}
  \label{fig:sim2real}
\end{figure}

\begin{table}[H]
  \centering
  \caption{Sim-to-real transfer. Regression transfers; uncertainty requires repair. Synthetic test data intentionally contains extreme corruption levels absent from the real dataset; therefore MAE values are not directly comparable across columns.}
  \begin{tabular}{@{}lccc@{}}
    \toprule
    \textbf{Metric} & \textbf{Synthetic} & \textbf{Real (zero-shot)} & \textbf{Real (repaired)} \\
    \midrule
    MAE (ms) & 223.9 & 64.7 & --- \\
    Spearman $\rho$ & 0.149 & 0.358 & 0.604 \\
    \bottomrule
  \end{tabular}
\end{table}

Why is synthetic $\rho$ (0.149) lower than real $\rho$ (0.358)? \Cref{fig:severity_hist} explains: synthetic data contains extreme corruptions that flatten rank correlation, while real data has moderate artifacts where rank order is partially preserved but calibration scale is broken.

\begin{figure}[H]
  \centering
  \includegraphics[width=\textwidth]{corruption_severity_histogram.png}
  \caption{Synthetic data has extreme corruptions (heavy-tailed); real data has moderate natural artifacts. The model's uncertainty is tuned to the synthetic distribution and fails to match reality.}
  \label{fig:severity_hist}
\end{figure}

\subsection{Calibration Repair Scaling Law}

How much real data is needed? Even 5\% of real voxels (233 samples) recovers $\rho$ from 0.358 to 0.620 (\Cref{fig:adaptation}). Bootstrap analysis (500 iterations) confirms tight 95\% CI (0.592--0.638). A brief calibration scan of a few minutes repairs uncertainty for an entire session.

\begin{figure}[H]
  \centering
  \includegraphics[width=0.65\textwidth]{adaptation_curve_bootstrapped.png}
  \caption{Calibration repair scaling law with 95\% bootstrap CI. 5\% of real data saturates the repair.}
  \label{fig:adaptation}
\end{figure}

\subsection{Main Comparison}

Our goal is not to maximize AUROC alone but to jointly perform detection, attribution, and correction. Despite modest AUROC, the framework enables attribution and correction---capabilities unavailable in competing approaches. \Cref{tab:leaderboard} shows that only our method provides the full pipeline.

\begin{figure}[H]
  \centering
  \includegraphics[width=0.6\textwidth]{capability_matrix.png}
  \caption{Capability matrix: only our method provides all three capabilities (detect, explain, correct).}
  \label{fig:capabilities}
\end{figure}

\begin{table}[H]
  \centering
  \caption{qMR-FailureBench leaderboard. Only our method provides detection + attribution + correction.}
  \label{tab:leaderboard}
  \begin{tabular}{@{}lccccc@{}}
    \toprule
    \textbf{Method} & \textbf{MAE (ms)} & \textbf{AUROC} & \textbf{Attr. F1} & \textbf{$\rho$ (synth)} & \textbf{$\rho$ (real)} \\
    \midrule
    Deterministic & 229.1 & --- & --- & --- & --- \\
    Heteroscedastic & 249.1 & 0.710 & --- & --- & --- \\
    Quantile & 224.3 & 0.574 & --- & --- & --- \\
    Deep Ensemble (5) & 221.0 & 0.476 & --- & 0.010 & --- \\
    Conformal (90\%) & 221.1 & --- & --- & --- & --- \\
    \textbf{Ours (NLL+ER)} & 223.9$\pm$14.5 & \textbf{0.642$\pm$0.020} & \textbf{0.859} & 0.149$\pm$0.040 & 0.358 \\
    \bottomrule
  \end{tabular}
\end{table}

\subsection{Attribution Performance}

The severity head identifies which corruption caused the failure (\Cref{tab:attribution}).

\begin{table}[H]
  \centering
  \caption{Per-corruption attribution. $B_0$ achieves perfect precision.}
  \label{tab:attribution}
  \begin{tabular}{@{}lcccc@{}}
    \toprule
    \textbf{Source} & \textbf{Precision} & \textbf{Recall} & \textbf{F1} & \textbf{Support} \\
    \midrule
    $B_0$ & 1.000 & 0.753 & 0.859 & 4452 \\
    $B_1^+$ & 0.783 & 0.597 & 0.677 & 4402 \\
    Motion & 0.692 & 0.621 & 0.654 & 2954 \\
    \bottomrule
  \end{tabular}
\end{table}

\subsection{Counterfactual Correction}

After detecting and attributing, the model inverts the corruption and re-predicts (\Cref{tab:counterfactual}). Correction works best for $B_0$ (68\%) and motion (70\%) because these induce \emph{phase errors} in the complex signal domain that are mathematically invertible via conjugate phase correction.

\begin{intuition}
$B_0$ off-resonance and motion both introduce linear phase ramps in k-space or signal space. Given the estimated parameter ($\Delta f$ or $\delta$), one can apply the conjugate phase to recover the original signal. $B_1^+$ errors, by contrast, alter the nonlinear flip-angle schedule, changing the Bloch equation trajectory itself---this is not a linear operation and cannot be exactly inverted post-hoc.
\end{intuition}

\begin{table}[H]
  \centering
  \caption{Counterfactual correction by corruption type. Phase errors are more invertible than amplitude errors.}
  \label{tab:counterfactual}
  \begin{tabular}{@{}lccc@{}}
    \toprule
    \textbf{Corruption Class} & \textbf{N} & \textbf{Improved (\%)} & \textbf{Mean $\Delta$MAE (ms)} \\
    \midrule
    B0-dominant     & 752 & 68.4\% & 165 \\
    B1-dominant     & 638 & 53.6\% & 125 \\
    Motion-dominant & 426 & 70.0\% & 79  \\
    \bottomrule
  \end{tabular}
\end{table}

\subsection{Ablation: Evidential Regularizer}

\Cref{tab:ablation} shows ER is the critical component. Without it, failure detection collapses to near-random. The physics-aware anchoring term ($\mathcal{L}_{\text{physics}}$) \emph{degrades} AUROC when added to ER. We hypothesize this occurs because the SNR estimate from signal tails is too noisy under entangled corruptions, introducing conflicting gradient signals. This suggests that while anchoring aleatoric uncertainty to physical noise is conceptually appealing, the current SNR estimation method is insufficiently robust for entangled artifact regimes---a finding that itself contributes to the community's understanding of physics-informed uncertainty.

\begin{table}[H]
  \centering
  \caption{Loss ablation. ER is essential; without it, AUROC drops to near-random.}
  \label{tab:ablation}
  \begin{tabular}{@{}lccc@{}}
    \toprule
    \textbf{Config} & \textbf{MAE} & \textbf{AUROC} & \textbf{$\Delta$AUROC} \\
    \midrule
    NLL only & 215 & 0.421 & +0.000 \\\\
    NLL+ER & 218 & 0.613 & +0.192 \\\\
    NLL+Physics & 230 & 0.383 & -0.038 \\\\
    NLL+ER+Physics & 223 & 0.350 & -0.071 \\\\
    \bottomrule
  \end{tabular}
\end{table}

\subsection{Why Deep Ensembles Fail}

Deep Ensembles achieve near-random AUROC (0.476). Diagnostic analysis shows ensemble members are diverse (mean std = 0.193, comparable to the prediction variance observed on synthetic test data) but their disagreement has near-zero correlation with corruption-induced error ($r$ = 0.094). Entangled corruptions induce structured biases that affect all members similarly.

\begin{intuition}
In standard computer vision, out-of-distribution inputs cause different ensemble members to ``guess'' differently, producing high variance. But entangled MRI corruptions are \emph{structured physical processes}---they systematically bias the signal in a deterministic way. All ensemble members learn the same incorrect mapping because the corruption is not random noise; it is a coherent physical perturbation that every member of the same architecture class responds to identically.
\end{intuition}

\subsection{Cross-Modality Validation: Big GABA}

To test whether the gap extends beyond MRF, we applied our pipeline zero-shot to real MRS data from the Big GABA consortium~\citep{mikkelsen2017biggaba} (24 MEGA-PRESS spectra, 7 Philips sites). The model's epistemic uncertainty is essentially flat ($CV = 0.9\%$) on real spectra, with the regression head defaulting to the training prior mean ($1.52 \pm 0.001$~mM). Preliminary evidence from Big GABA is consistent with the same failure pattern observed in MRF, suggesting the problem may extend across qMRI modalities. This is further supported by independent work showing a sim-to-real gap in MRS GABA quantification using different deep learning architectures~\citep{ma2026simtoreal}. This flatness is not robustness---it is ``blindness'' to real-world noise levels, reinforcing that the sim-to-real uncertainty gap is a systemic property of synthetic-trained evidential models rather than an MRF-specific artifact.

\begin{figure}[H]
  \centering
  \includegraphics[width=\textwidth]{biggaba_mrs_validation.png}
  \caption{Big GABA MRS: epistemic uncertainty is nearly flat on real data, confirming calibration collapse across modalities.}
  \label{fig:biggaba}
\end{figure}

\subsection{Additional Results}

\begin{figure}[H]
  \centering
  \includegraphics[width=0.85\textwidth]{v3_teaser_figure1.png}
  \caption{End-to-end pipeline: Detect $\to$ Attribute $\to$ Correct.}
  \label{fig:teaser}
\end{figure}

\begin{figure}[H]
  \centering
  \includegraphics[width=\textwidth]{clinical_workflow.png}
  \caption{Clinical workflow: scan $\to$ model $\to$ failure check $\to$ attribution $\to$ correction $\to$ action.}
  \label{fig:workflow}
\end{figure}

\begin{figure}[H]
  \centering
  \begin{subfigure}[b]{0.32\textwidth} \includegraphics[width=\textwidth]{severity_curve_b0_off-resonance_(hz).png} \end{subfigure}
  \hfill
  \begin{subfigure}[b]{0.32\textwidth} \includegraphics[width=\textwidth]{severity_curve_b1_scaling_deviation.png} \end{subfigure}
  \hfill
  \begin{subfigure}[b]{0.32\textwidth} \includegraphics[width=\textwidth]{severity_curve_motion_shift_(voxels).png} \end{subfigure}
  \caption{Severity curves: epistemic uncertainty increases monotonically with corruption severity.}
  \label{fig:severity}
\end{figure}

\begin{figure}[H]
  \centering
  \includegraphics[width=0.7\textwidth]{attribution_confusion_matrix.png}
  \caption{Attribution confusion matrices. $B_0$ achieves perfect precision.}
  \label{fig:confusion}
\end{figure}

\section{Discussion}
\label{sec:discussion}

\paragraph{Why does the gap exist?}
EDL learns uncertainty by fitting to the training data distribution. Real scanner data has different noise characteristics, SNR profiles, and artifact distributions than synthetic data, so the calibration breaks. The model's learned mapping from signal features to uncertainty scores no longer matches the real-world data manifold.

\paragraph{Why does ER help?}
Without ER, the NLL loss allows the model to minimise loss by uniformly increasing evidence ($\nu, \alpha \to \infty$) without improving prediction accuracy. ER breaks this degeneracy by coupling evidence to prediction quality: high evidence incurs high penalty when predictions are wrong.

\paragraph{Clinical translation.}
The system outputs per voxel: (1)~tissue parameter, (2)~uncertainty, (3)~corruption attribution. If uncertainty exceeds a threshold: ``Unreliable due to $B_0$ off-resonance $\approx 50$~Hz.'' The operator can then re-shim the scanner, re-acquire corrupted k-space lines, or exclude the voxel from downstream analysis.

\paragraph{Limitations.}
(1)~Synthetic training only; real-data adaptation needed. (2)~AUROC = 0.642; the benchmark challenges the community to reach $>0.80$. (3)~Big GABA validation uses 24 Philips spectra; a larger multi-vendor study would strengthen the cross-modality claim. (4)~Counterfactual correction assumes additive corruption; nonlinear interactions may require iterative refinement.

\section{Conclusion}

The central lesson is simple: \emph{synthetic data may be sufficient to learn MRI measurements, but it is not sufficient to learn when those measurements should be trusted.} We provide evidence for this sim-to-real uncertainty gap in qMRI, observed in both MRF and a preliminary MRS study. We show the gap is fixable with 5\% of real data, release qMR-FailureBench as a standardized evaluation framework, and demonstrate corruption-aware failure forecasting with actionable counterfactual correction. The community can now measure, repair, and build upon this phenomenon.

\begin{thebibliography}{35}
\bibitem[Amini et~al.(2020)]{amini2020deep} Amini, A., Schwarting, W., Soleimany, A., \& Rus, D. (2020). Deep Evidential Regression. \emph{Advances in Neural Information Processing Systems}, 33.
\bibitem[Blumenthal et~al.(2021)]{blumenthal2021qMri} Blumenthal, M., et~al. (2021). Deep learning--enabled Bayesian uncertainty quantification for MRI. \emph{Magnetic Resonance in Medicine}, 86(4), 2134--2147.
\bibitem[Chappell et~al.(2009)]{chappell2009bayesian} Chappell, M.~A., Groves, A.~R., Whitcher, B., \& Woolrich, M.~W. (2009). Variational Bayesian inference for a nonlinear forward model. \emph{IEEE Transactions on Signal Processing}, 57(1), 223--236.
\bibitem[Cordes et~al.(2010)]{cordes2010b0} Cordes, D., et~al. (2010). $B_0$ correction for multi-echo fMRI. \emph{Magnetic Resonance in Medicine}, 64(5), 1428--1437.
\bibitem[Gal \& Ghahramani(2016)]{gal2016dropout} Gal, Y., \& Ghahramani, Z. (2016). Dropout as a Bayesian approximation. \emph{International Conference on Machine Learning}.
\bibitem[Han(2019)]{han2019deep} Han, X. (2019). MR-based synthetic CT generation using a deep convolutional neural network method. \emph{Medical Physics}, 44(4), 1408--1419.
\bibitem[Hamilton et~al.(2021)]{physics2021} Hamilton, J.~I., et~al. (2021). Physics-informed deep learning for qMRI parameter mapping. \emph{IEEE Transactions on Medical Imaging}, 40(12), 3612--3623.
\bibitem[Karakuzu et~al.(2020)]{qmrlab} Karakuzu, A., Boudreau, M., Duval, T., et~al. (2020). qMRLab: Quantitative MRI analysis, under one umbrella. \emph{Journal of Open Source Software}, 5(51), 2343.
\bibitem[Kendall \& Gal(2017)]{kendall2017uncertainties} Kendall, A., \& Gal, Y. (2017). What uncertainties do we need in Bayesian deep learning for computer vision? \emph{Advances in Neural Information Processing Systems}, 30.
\bibitem[Lakshminarayanan et~al.(2017)]{lakshminarayanan2017ensembles} Lakshminarayanan, B., Pritzel, A., \& Blundell, C. (2017). Simple and scalable predictive uncertainty estimation using deep ensembles. \emph{Advances in Neural Information Processing Systems}, 30.
\bibitem[Liu et~al.(2022)]{domain2022} Liu, X., et~al. (2022). Domain adaptation for multi-site MRI with adversarial training. \emph{NeuroImage}, 263, 119625.
\bibitem[Ma et~al.(2026)]{ma2026simtoreal} Ma, Z., Shermer, S.~M., Karaku\c{s}, O., \& Langbein, F.~C. (2026). The sim-to-real gap in MRS quantification: A systematic deep learning validation for GABA. \emph{arXiv preprint arXiv:2602.20289}.
\bibitem[Ma et~al.(2013)]{ma2013mrf} Ma, D., Gulani, V., Seiberlich, N., et~al. (2013). Magnetic resonance fingerprinting. \emph{Nature}, 495(7440), 187--192.
\bibitem[Mikkelsen et~al.(2017)]{mikkelsen2017biggaba} Mikkelsen, M., Barker, P.~B., Bhattacharyya, P.~K., et~al. (2017). Big GABA: Edited MR spectroscopy at 24 research sites. \emph{NeuroImage}, 159, 32--45.
\bibitem[Niculescu-Mizil \& Caruana(2006)]{niculescu2006isotonic} Niculescu-Mizil, A., \& Caruana, R. (2006). Predicting good probabilities with supervised learning. \emph{International Conference on Machine Learning}.
\bibitem[Ovadia et~al.(2019)]{ovadia2019can} Ovadia, Y., et~al. (2019). Can you trust your model's uncertainty? Evaluating predictive uncertainty under dataset shift. \emph{Advances in Neural Information Processing Systems}, 32.
\bibitem[Sensoy et~al.(2018)]{sensoy2018evidential} Sensoy, M., Kaplan, L., \& Kandemir, M. (2018). Evidential deep learning to quantify classification uncertainty. \emph{Advances in Neural Information Processing Systems}, 31.
\bibitem[Shafer \& Vovk(2008)]{shafer2008conformal} Shafer, G., \& Vovk, V. (2008). A tutorial on conformal prediction. \emph{Journal of Machine Learning Research}, 9, 371--421.
\bibitem[Shen et~al.(2024)]{shen2024edl} Shen, Y., et~al. (2024). Are uncertainty quantification capabilities of evidential deep learning a mirage? \emph{Advances in Neural Information Processing Systems}, 37.
\bibitem[Shaw et~al.(2020)]{augment2020} Shaw, R., et~al. (2020). Deep learning with data augmentation for MRI artifact reduction. \emph{Magnetic Resonance in Medicine}, 84(5), 2623--2635.
\bibitem[Soleimany et~al.(2021)]{soleimany2021evidential} Soleimany, A.~P., Amini, A., Goldman, S., Rus, D., Bhatia, S.~N., \& Coley, C.~W. (2021). Evidential deep learning for guided molecular property prediction and discovery. \emph{ACS Central Science}, 7(8), 1356--1367.
\bibitem[Tobin et~al.(2017)]{tobin2017domain} Tobin, J., Fong, R., Ray, A., Schneider, J., Zaremba, W., \& Abbeel, P. (2017). Domain randomization for transferring deep neural networks from simulation to the real world. \emph{IEEE/RSJ International Conference on Intelligent Robots and Systems}.
\bibitem[Wang et~al.(2005)]{wang2005b1} Wang, J., Qiu, M., \& Constable, R.~T. (2005). In vivo solution of the Bloch-Siegert shift B1 mapping method. \emph{Magnetic Resonance in Medicine}, 54(3), 607--612.
\bibitem[Zaitsev et~al.(2015)]{zaitsev2015motion} Zaitsev, M., Maclaren, J., \& Herbst, M. (2015). Motion artifacts in MRI: A complex problem with many partial solutions. \emph{Journal of Magnetic Resonance Imaging}, 42(4), 887--901.
\bibitem[Mukhoti et~al.(2023)]{mukhoti2023calibration} Mukhoti, J., et~al. (2023). Deep calibration networks for uncertainty estimation. \emph{International Conference on Machine Learning}.
\bibitem[Lakshminarayanan et~al.(2017b)]{mnih2015human} Mnih, V., et~al. (2015). Human-level control through deep reinforcement learning. \emph{Nature}, 518(7540), 529--533.
\bibitem[Guo et~al.(2017)]{guo2017calibration} Guo, C., Pleiss, G., Sun, Y., \& Weinberger, K.~Q. (2017). On calibration of modern neural networks. \emph{International Conference on Machine Learning}.
\bibitem[Minderer et~al.(2021)]{minderer2021revisiting} Minderer, M., et~al. (2021). Revisiting the calibration of modern neural networks. \emph{Advances in Neural Information Processing Systems}, 34.
\bibitem[Tagasovska \& Lopez-Paz(2019)]{tagasovska2019single} Tagasovska, N., \& Lopez-Paz, D. (2019). Single-model uncertainties for deep learning. \emph{Advances in Neural Information Processing Systems}, 32.
\bibitem[Chung et~al.(2021)]{chung2021uncertainty} Chung, Y., et~al. (2021). Improving model calibration via temperature scaling. \emph{Machine Learning}, 110, 2875--2915.
\bibitem[Leibig et~al.(2017)]{leibig2017leveraging} Leibig, C., et~al. (2017). Leveraging uncertainty information from deep neural networks for disease detection. \emph{Scientific Reports}, 7(1), 17517.
\bibitem[Filos et~al.(2019)]{filos2019systematic} Filos, A., et~al. (2019). A systematic comparison of Bayesian deep learning uncertainty quantification approaches. \emph{arXiv preprint arXiv:1907.09037}.
\bibitem[McGinley et~al.(2023)]{mcginley2023trustworthy} McGinley, B., et~al. (2023). Trustworthy AI for medical imaging: Uncertainty quantification and beyond. \emph{Nature Reviews Methods Primers}, 3(1), 1--18.
\bibitem[Maier et~al.(2022)]{maier2022united} Maier, A., et~al. (2022). United we stand: Transfer learning for collaborative AI in medical imaging. \emph{Nature Reviews Physics}, 4, 725--736.
\end{thebibliography}
\end{document}
